%------------------------------%
%% ✎ Dylan (V1) %%%%%%%%% ✅ %%
%% ✎ Alain (V2) %%%%%%%%% ✅ %%
%% ✎ Dylan (V3) %%%%%%%%% ✅ %%
%------------------------------%

\afterpage{%
%\afterpage{%

    % Arrière-plan partie I
    \AddToShipoutPictureBG*{%
\includegraphics[width=\paperwidth,height=\paperheight]{src/Figures/Arriere_plan/Arriere_plan_Part_1.jpg}
    }

% Rectangle
\AddToShipoutPictureBG*{
  \begin{tikzpicture}[remember picture,overlay]
    \node[fill=white, opacity=0.75, text width=\paperwidth, minimum height=12cm, anchor=north] 
    at ([yshift=-6.7cm]current page.north) {};
  \end{tikzpicture}
}

% Source
\AddToShipoutPictureFG*{
  \AtPageLowerRight{
    \raisebox{1cm}{
      \hspace{16cm}
      \begin{tikzpicture}
        \node[fill=white, rounded corners=5pt, inner sep=5pt, align=center] {
          \tiny{Photography: \textcolor{blue}{Dylan Moinse (2021)}}
        };
      \end{tikzpicture}
    }
  }
}
}%}

\renewcommand{\thefigure}{\thechapter.\arabic{figure}} % Numérotation standard
\renewcommand{\thetable}{\thechapter.\arabic{table}}
\setcounter{figure}{0}
\setcounter{table}{0}

\needspace{1\baselineskip} % Réserve de l'espace
\part{Theoretical and Methodological Frameworks for Integrating Light Individual Mobility in the Context of Rail-Oriented Urbanism
    \label{part1:titre}
    }
    \markboth{Part~I: Theoretical and Methodological Frameworks of the Research}{}
    \markright{Part~I: Theoretical and Methodological Frameworks of the Research}{}

% Introduction of Part I
\cleardoublepage
\section*{Introduction of Part~I
    \label{part1:introduction}
    }
    \addcontentsline{toc}{chapter}{Introduction of Part~I}

    % Introduction
\lettrine[lines=3, findent=8pt, nindent=0pt]{\lettrinefont T}{he} integration of light individual mobility in station districts and its role in reinterpreting \acrfull{TOD} is part of a broader reflection on the interactions between urbanism, territorial configurations, transport infrastructures, and mobility practices. Before embarking on the construction of empirical data and its analysis, it is essential to revisit the theoretical foundations, conceptual developments, and research methods that will structure this study. This first part thus aims to establish a theoretical and methodological framework that enables understanding the challenges of intermodal accessibility and exploring the ways in which light individual mobility can be integrated into urban planning and transport strategies. It is based on the articulation of three dimensions: (i) a conceptual review of urban models and their adaptation to contemporary mobility dynamics; (ii) a state-of-the-art review of existing research on intermodality and the complementarity between individual mobility and public transport; (iii) a methodological framework detailing the approaches, tools, and analytical techniques used to address these transformations. By gathering these elements, this first part anchors the reflection within a structured framework, justifying the interest of examining \acrshort{TOD} not as a static urban model, but as an evolving strategy capable of integrating light individual mobility to maximize the expected benefits.%%Translated%%

    % Chapter 1
\textsl{Shedding Light on the Conceptual Evolution of Rail-Based Urbanism in View of Its Adjustments} (\hyperref[objectif-1]{Objective~\(O_1\)}, page~\pageref{objectif-1}). \hyperref[chap1:titre]{The first chapter} (page~\pageref{chap1:titre}) traces the evolution of \acrshort{TOD} as a model for structuring regional urban development around public transport networks. It analyzes its core principles, including densification around transport hubs, functional mix, and the treatment of public spaces as tools for creating territories less dependent on automobiles. Despite its ambitions, \acrshort{TOD} presents several structural limitations, particularly in terms of accessibility. While its effectiveness is evident in dense metropolitan contexts, where public transport infrastructures are well integrated into the urban fabric, it proves less effective in suburban areas, where dependence on cars remains prevalent. The main cause lies in the exclusion of intermodal solutions tailored to the \Commas{first and last mile} problem. This situation thus limits the expected positive impacts of \acrshort{TOD} in terms of modal shift and transport network optimization. It is in this context that light individual mobility—including bicycles, scooters, and other micromobility devices, which will also require a reconsideration of their evolution—appears as a strategic lever to complement and update \acrshort{TOD}. Its integration into urban planning and transport policies could help fill the gaps of the traditional model by ensuring a more efficient and alternative mobility system aligned with emerging aspirations. However, this complementarity remains underexplored in academic research and insufficiently implemented in public policies. Once the interest of \acrshort{TOD} as a theoretical framework to address the need for reducing car usage has been demonstrated, its adaptability in areas historically designed for automobiles should be questioned. This reflection thus leads to reviewing the existing literature on this issue and drawing the necessary lessons to guide our own investigation.%%Translated%%

    % Chapter 2
\textsl{Synthesizing Research Practices and Current Knowledge on Public Transport-Oriented Urbanism Supported by Light Individual Mobility} (\hyperref[objectif-2]{Objective~\(O_2\)}, page~\pageref{objectif-2}). \hyperref[chap2:titre]{The second chapter} (page~\pageref{chap2:titre}) constitutes a \acrfull{SLR} designed to better understand the potential synergies between \acrshort{TOD} and light individual mobility. In this context, we use a scientometric approach to identify major research trends related to this topic and to analyze the evolution of associated concepts. This chapter aims to survey and map the works that, implicitly or explicitly, involve the contours of a \acrshort{B-TOD} or \acrshort{M-TOD}. It focuses not only on the geographical and temporal dynamics of publications and their semantics but also on the methods employed in these studies to understand their analytical approaches and underlying methodological frameworks. The analysis of these developments allows for a refined understanding of intermodal mobility practices, incorporating a more nuanced reading of the uses and spatial determinants of intermodal accessibility. Finally, this critical literature review highlights the unresolved challenges, both conceptually and in terms of practical implementation. It raises several key questions for the analyses to come, identifying scientific and methodological gaps that will guide the empirical work of this research. This state-of-the-art review thus serves as a prerequisite for the empirical investigation, building on the gaps identified in the literature to propose a renewed approach capable of better addressing the challenges of \acrshort{M-TOD}.%%Translated%%

    % Chapter 3
\textsl{Developing a Methodology Integrating Complementary Approaches to Achieve a More In-Depth Study and Extract Additional Results} (\hyperref[objectif-3]{Objective~\(O_3\)}, page~\pageref{objectif-3}). \hyperref[chap3:titre]{The third chapter} (page~\pageref{chap3:titre}) addresses the empirical questions raised in the previous chapters. To this end, we have designed and adapted an original methodological framework. This chapter thus presents the methodological approach adopted in this thesis. First, it defines and justifies the choice of the study area, both in terms of the objects analyzed and the geographic area concerned. The construction of our empirical data involves the use of mixed approaches, combining the exploitation of databases, a survey that is both quantitative and qualitative, and spatial modeling. The methodology described allows for the articulation of three complementary levels of analysis: a behavioral level, focusing on modal choices and mobility strategies; an infrastructural level, exploring the physical and organizational conditions of accessibility; and a territorial level, evaluating the potential and actual impact of light individual mobility on station districts. Furthermore, this chapter details the methods used to process the collected data, namely the tools and analytical techniques employed at the interface between geography, sociology, economics, and data science.%%Translated%%